\documentclass[10pt,a4paper]{amsart}
\usepackage[margin=19mm]{geometry}
\usepackage{amsmath,amssymb,mathtools}
\usepackage[scr=rsfs,cal=euler]{mathalfa}
\usepackage{xfrac}
\usepackage[unicode,hidelinks,bookmarks=false]{hyperref}
\newcommand{\ie}{i.e.\ }
\newcommand{\cf}{cf.\ }
\newcommand{\resp}{respectively\ }
\newcommand{\Subsection}{\subsection}
\providecommand{\nobreakdash}{\nobreak\hbox{-}}
\setlength{\emergencystretch}{2em}
\multlinegap0pt
\usepackage{bm}
\usepackage{bbm}
\usepackage{dsfont}
\usepackage{xspace}
\usepackage{cancel}
\usepackage{mathtools}
\usepackage{amsthm}
\makeatletter
\@ifundefined{theorem}{\newtheorem{theorem}{Theorem}}{}
\makeatother
\title[Measure and integration]{Measure and integration}
\author{Teo Banica}
\date{Source: arXiv:2506.15534v1 (CC BY 4.0)}
\begin{document}
\maketitle

\noindent\emph{Source and licence.} Measure and integration. Version arXiv:2506.15534v1 (CC BY 4.0), distributed under the Creative Commons Attribution 4.0 International licence, as recorded in the arXiv licence metadata for this exact version and shown on the abstract page https://arxiv.org/abs/2506.15534v1. Full rights notice: licenses/arxiv-2506.15534-CC-BY-4.0.txt.\par\smallskip

\noindent\textit{Editorial scope.} Counted unit: the Theorem whose source label is \texttt{None} in the article (source0.tex lines 5481--5493, source label \texttt{None}), delivered together with the article's own complete original proof of it (source0.tex lines 5495--5834, one \texttt{proof} environment of 340 lines). No mathematics is written, repaired or re-derived: statement and proof are the article's own text line for line. Required context retained uncounted: none. Every definition, display and label of those lines is retained unchanged. Omitted from this package and not redistributed: 1--5480, 5494--5494, 5835--17104. Nothing inside a retained excerpt is omitted or altered: every retained body is delivered exactly as the article's lines are written, so no omission receipt of any kind accompanies this package. Citations: the retained excerpts carry no citation command at all, so no bibliography entry of the article is transcribed and no reference is redistributed. Reference adaptation: none was needed and none was made; every reference inside the delivered unit resolves inside it, so no unresolved reference or citation is produced. Printed slips kept literal: no printed word, symbol, hypothesis or number of the counted statement or of its proof was altered. Layout and class: the article's own class is \texttt{amsbook} with packages including amssymb, xy, hyperref; the wrapper is the lane's pinned \texttt{amsart} editorial wrapper at 10pt on A4, which restates the theorem-like environments the retained text needs and adds the article's own macro definitions that the retained text needs (none), with extra wrapper packages (bm, bbm, dsfont, xspace, cancel, mathtools). The delivered numbering is the unit's own. Assets: no retained span contains a figure environment, a graphics-inclusion command, a PDF-page inclusion, a table or a TikZ picture; no image file is copied into this package, the package declares no asset dependency and no image file is redistributed with it. Licence: version arXiv:2506.15534v1 of this article is distributed under the Creative Commons Attribution 4.0 International licence, as recorded in the arXiv licence metadata for this exact version and shown on the abstract page https://arxiv.org/abs/2506.15534v1. A full rights notice is delivered at licenses/arxiv-2506.15534-CC-BY-4.0.txt. Characters: every retained excerpt is pure ASCII, so the delivered engine, XeLaTeX with its default Latin Modern text font, reports no missing character.
\begin{theorem}[Riesz]
Any positive functional $I:C_c(X)\to\mathbb R$ comes by integrating with respect to a certain measure on $X$,
$$I(f)=\int_Xf(x)d\mu(x)$$
which has the following additional properties,
\begin{enumerate}
\item $\mu(K)<\infty$, for any $K\subset X$ compact.

\item $\mu(E)=\inf\{\mu(U)|E\subset U\ {\rm open}\}$, for any $E$ measurable.

\item $\mu(E)=\sup\{\mu(K)|K\subset E\ {\rm compact}\}$, for any $E$ open, or of finite measure.
\end{enumerate}
and which is unique with these properties, modulo the null sets.
\end{theorem}

\begin{proof}
This is something quite long and tricky, the idea being as follows:

\medskip

(1) Let us first prove the uniqueness. Assuming that our measure $\mu$ produces $I$, and has the various properties in the statement, it is clear that $\mu$ is uniquely determined by its values on the compact sets $K\subset X$. Thus, we must show that given two measures $\mu_1,\mu_2$ as in the statement, and a compact set $K\subset X$, we have:
$$\mu_1(K)=\mu_2(K)$$

For this purpose, let us pick $\varepsilon>0$. By using the properties (1,3) in the statement, for the measure $\mu_2$, we can find $K\subset U$ open such that:
$$\mu_2(U)<\mu_2(K)+\varepsilon$$

Now by using the Urysohn lemma for the inclusion $K\subset U$, we obtain a certain continuous, compactly supported function $f$, such that:
$$\chi_K\leq f\leq\chi_U$$

But, with this choice of $f$, we have the following computation:
\begin{eqnarray*}
\mu_1(K)
&\leq&\int_Xf(x)d\mu_1(x)\\
&=&\int_Xf(x)d\mu_2(x)\\
&\leq&\mu_2(U)\\
&<&\mu_2(K)+\varepsilon
\end{eqnarray*}

Thus $\mu_1(K)\leq\mu_2(K)$, and by interchanging $\mu_1,\mu_2$ we have the reverse inequality as well, so we obtain, as desired, $\mu_1(K)=\mu_2(K)$, proving the uniqueness statement.

\medskip

(2) With this done, let us get now to the main part, and discuss the construction of the algebra $M$, that we will choose to be saturated in the sense of Theorem 6.5, and of the measure $\mu$. What we have as data is a positive functional, as follows:
$$I:C_c(X)\to\mathbb R$$

We must first measure the Borel sets $E\subset X$. And here, to start with, in order to measure the open sets $U\subset X$, the formula is straightforward, namely:
$$\mu(U)=\sup\left\{I(f)\Big| f\leq\chi_U\right\}$$

Now observe that with this definition in hand, for the open sets, we have:
$$U_1\subset U_2\implies\mu(U_1)\leq\mu(U_2)$$

We deduce that the following happens, for the open sets:
$$\mu(E)=\inf\left\{\mu(U)\Big| E\subset U\ {\rm open}\right\}$$

But this can serve as the definition of $\mu$, for all the subsets $E\subset X$.

\medskip

(3) Regarding now the measurable sets, let us first consider the class $N$ of subsets $E\subset X$ which are of finite measure, $\mu(E)<\infty$, and satisfy the following condition:
$$\mu(E)=\sup\left\{\mu(K)\Big| K\subset E\ {\rm compact}\right\}$$

Then, we can define the class of measurable sets $M\subset P(X)$ as follows:
$$M=\left\{E\subset X\Big|E\cap K\in N,\forall K\subset X\ {\rm compact}\right\}$$

Summarizing, done with definitions, and it remains now to prove that $M$ is an algebra, and that $\mu$ is a measure on it, along with the other things claimed in the statement.

\medskip

(4) In order to prove that $\mu$ is indeed a measure, our first claim is that we have the following inequality, for any two open sets $U_1,U_2\subset X$:
$$\mu(U_1\cup U_2)\leq\mu(U_1)+\mu(U_2)$$ 

In order to prove this claim, choose an arbitrary function $g:X\to[0,1]$, supported on $U_1\cup U_2$. By using Theorem 6.11 we obtain a certain partition of unity, $f_1+f_2=1$ on the support of $g$, and it follows that we have the following equality:
$$g=f_1g+f_2g$$

Now by applying our integration functional $I$, we obtain from this:
\begin{eqnarray*}
I(g)
&=&I(f_1g+f_2g)\\
&=&I(f_1g)+I(f_2g)\\
&\leq&\mu(U_1)+\mu(U_2)
\end{eqnarray*}

Now since this holds for any $g:X\to[0,1]$ as above, this proves our claim.

\medskip

(5) Our second claim is that we have in fact, more generally, the following inequality, valid for any family of subsets $E_1,E_2,E_3,\ldots$ of our space $X$:
$$\mu\left(\bigcup_{i=1}^\infty E_i\right)\leq\sum_{i=1}^\infty\mu(E_i)$$

Indeed, this inequality trivially holds if $\mu(E_i)=\infty$ for some $i$, so we can assume $\mu(E_i)<\infty$ for any $i$. Now pick $\varepsilon>0$, and choose open sets $U_i\supset E_i$ such that:
$$\mu(U_i)<\mu(E_i)+\frac{\varepsilon}{2^i}$$

Consider also the union of these open sets $U_i$, chosen as above:
$$U=\bigcup_{i=1}^\infty U_i$$

Now pick an arbitrary function $f:X\to[0,1]$, supported on this open set $U$. Since $f$ has compact support, we can find a certain integer $n\in\mathbb N$ such that:
$$supp(f)\subset U_1\cup\ldots\cup U_n$$

Now by using what we found in (4), recursively, we obtain:
\begin{eqnarray*}
I(f)
&\leq&\mu(U_1\cup\ldots\cup U_n)\\
&\leq&\mu(U_1)+\ldots+\mu(U_n)\\
&\leq&\sum_{i=1}^n\mu(E_i)+\varepsilon
\end{eqnarray*}

Since this holds for any $f:X\to[0,1]$ as above, we deduce that we have:
$$\mu\left(\bigcup_{i=1}^\infty E_i\right)
\leq\mu(U)
\leq\sum_{i=1}^\infty\mu(E_i)+\varepsilon$$

Now since the number $\varepsilon>0$ was arbitrary, this proves our claim.

\medskip

(6) Our next claim, which proves in the assertion (1) in the theorem, is that any compact set $K\subset X$ is measurable, and that we have:
$$\mu(K)<\infty$$

In order to prove this claim, consider an arbitrary continuous function $f:X\to[0,1]$ satisfying $f(x)=1$ on $K$, then pick $\alpha<1$, and consider the following open set:
$$U_\alpha=\left\{x\in X\Big|f(x)>\alpha\right\}$$

We have then $K\subset U_\alpha$, and so $\alpha g\leq f$, whenever a continuous function $g:X\to[0,1]$ is supported by $U_\alpha$. By using this, we obtain the following estimate:
\begin{eqnarray*}
\mu(K)
&\leq&
\mu(U_\alpha)\\
&=&\sup\left\{I(g)\Big|g:X\to[0,1],\,supp(g)\subset U_\alpha\right\}\\
&\leq&\alpha^{-1}I(f)
\end{eqnarray*}

Now with $\alpha\to 1$, we obtain from this the following estimate:
$$\mu(K)\leq I(f)$$

Thus $\mu(K)<\infty$, as claimed, and the fact that $K$ is measurable is clear too.

\medskip

(7) Our next claim, which improves what we found in (6), is that for any compact set $K\subset X$ we have in fact the following estimate:
$$\mu(K)=\inf\left\{I(f)\Big|f:X\to[0,1],\,f(x)=1\ {\rm on}\ K\right\}$$

In order to prove this, let us go back to the proof of (6). We know from there that for any continuous function $f:X\to[0,1]$ satisfying $f(x)=1$ on $K$, as above, we have:
$$\mu(K)\leq I(f)$$

Now pick $\varepsilon>0$, and choose an open set $U\supset K$ satisfying:
$$\mu(U)<\mu(K)+\varepsilon$$

By using Lemma 6.10, we can find a compactly supported function $f$ such that:
$$\chi_K\leq f\leq\chi_U$$ 

But this gives the following estimate, using the above inequality $\mu(K)\leq I(f)$:
\begin{eqnarray*}
I(f)
&\leq&\mu(V)\\
&<&\mu(K)+\varepsilon\\
&\leq&I(f)+\varepsilon
\end{eqnarray*}

Now since the number $\varepsilon>0$ was arbitrary, this proves our claim.

\medskip

(8) Our claim now is that any open set $U\subset X$ satisfies the following condition, that we used in (3) in order to define the measurable sets:
$$\mu(U)=\sup\left\{\mu(K)\Big| K\subset U\ {\rm compact}\right\}$$

In order to prove this, pick an arbitrary number $\alpha<\mu(U)$. We can then find a continuous function $f:X\to[0,1]$ supported on $U$, such that:
$$I(f)>\alpha$$

Consider now the compact set $K=supp(f)$. Then for any open set $V$ containing $K$ we have $I(f)\leq\mu(V)$, and we deduce from this that we have:
$$I(f)\leq\mu(K)$$

In other words, given $\alpha<\mu(U)$, we have found a compact set $K$ such that:
$$\mu(K)\geq\alpha$$

Now since our number $\alpha<\mu(U)$ was arbitrary, this proves our claim.

\medskip

(9) Getting now towards the proof of the additivity of our measure $\mu$, let us first prove that for any two disjoint compact sets $K_1,K_2\subset X$, we have:
$$\mu(K_1\cup K_2)=\mu(K_1)+\mu(K_2)$$

For this purpose, pick an arbitrary $\varepsilon>0$. By using Lemma 6.10, we can find a compactly supported continuous function $f:X\to[0,1]$ satisfying:
$$f(x)=\begin{cases}
1&{\rm on}\ K_1\\
0&{\rm on}\ K_2
\end{cases}$$

On the other hand, by using (7) we can find a compactly supported continuous function $g:X\to[0,1]$ satisfying $g(x)=1$ on $K_1\cup K_2$, such that:
$$I(g)\leq\mu(K_1\cup K_2)+\varepsilon$$

By using this, and the linearity of $I$, we have the following estimate:
\begin{eqnarray*}
\mu(K_1)+\mu(K_2)
&\leq&I(fg)+I(g-fg)\\
&=&I(g)\\
&<&\mu(K_1\cup K_2)+\varepsilon
\end{eqnarray*}

Now since our number $\varepsilon>0$ was arbitrary, this proves our claim, via (5).

\medskip

(10) We are now in position of making a key verification, that of the additivity property of our measure $\mu$. To be more precise, our claim is that we have the following equality, valid for any family of pairwise disjoint subsets $E_1,E_2,E_3,\ldots$ of our space $X$:
$$\mu\left(\bigcup_{i=1}^\infty E_i\right)=\sum_{i=1}^\infty\mu(E_i)$$

In order to prove this, consider the union on the left, namely:
$$E=\bigcup_{i=1}^\infty E_i$$

In the case $\mu(E)=\infty$ our claim is proved by (5). So, assume $\mu(E)<\infty$, and pick an arbitrary number $\varepsilon>0$. We can then find compact sets $K_i\subset E_i$ such that:
$$\mu(K_i)>\mu(E_i)-\frac{\varepsilon}{2^i}$$

Now consider the following unions, which are compact sets too:
$$H_n=K_1\cup\ldots\cup K_n$$

By using the additivity formula that we found in (9), recursively, we obtain:
$$\mu(E)
\geq\mu(H_n)
=\sum_{i=1}^n\mu(K_i)
>\sum_{i=1}^n\mu(E_i)-\varepsilon$$

Now since our number $\varepsilon>0$ was arbitrary, this proves our claim.

\medskip

(11) Thus, additivity property proved, and with this in hand, the other assertions are quite standard. In order to prove these, our first claim is that for any measurable set $E$ and any $\varepsilon>0$, there is a compact set $K$, and an open set $U$, satisfying:
$$K\subset E\subset U\quad,\quad \mu(U-K)<\varepsilon$$

Indeed, in order to prove this claim, the first remark is that we can find indeed a compact set $K$ and an open set $U$, satisfying the following conditions:
$$K\subset E\subset U\quad,\quad \mu(U)-\frac{\varepsilon}{2}<\mu(E)<\mu(K)+\frac{\varepsilon}{2}$$

Since the set $U-K$ is open, by using (8) and then (10) we obtain:
$$\mu(U)=\mu(K)+\mu(U-K)$$

But with this, we can establish our inequality, as follows:
$$\mu(U-K)
=\mu(U)-\mu(K)<\varepsilon$$

(12) Our next claim is that if $A,B$ are measurable, then so are the following sets:
$$A-B\quad,\quad A\cup B\quad,\quad A\cap B$$

In order to prove this, pick an arbitrary number $\varepsilon>0$. By using (11) we can find certain compact sets $K,H$, and certain open sets $U,V$, satisfying:
$$K\subset A\subset U\quad,\quad \mu(U-K)<\varepsilon$$
$$H\subset B\subset V\quad,\quad \mu(V-H)<\varepsilon$$

Now observe that we have the following inclusion:
$$A-B
\subset U-H
\subset(U-K)\cup(K-V)\cup(V-H)$$

By using now (5), we obtain from this the following estimate:
\begin{eqnarray*}
\mu(A-B)
&\leq&\mu(U-K)+\mu(K-V)+\mu(V-H)\\
&<&\varepsilon+\mu(K-V)+\varepsilon\\
&=&\mu(K-V)+2\varepsilon
\end{eqnarray*}

Now since $K-V$ is a compact subset of $A-B$, we conclude that $A-B$ is measurable, as desired. Regarding now $A\cup B$, we can use here the following formula:
$$A\cup B=(A-B)\cup B$$

Indeed, by using (10) we obtain from this that $A\cup B$ is measurable as well. Finally, regarding  $A\cap B$, we can use here the following formula:
$$A\cap B=A-(A-B)$$

Thus $A\cap B$ is measurable as well, and this finishes the proof of our claim.

\medskip

(13) Our claim now is that the set $M$ constructed in (3) is indeed an algebra, which contains all Borel sets. In order to prove this, consider an arbitrary compact set $K\subset X$. Given $A\in M$, we can write $A^c\cap K$ as a difference of two sets in $M$, as follows:
$$A^c\cap K=K-(A-K)$$

Thus $A^c\cap K\in M$, and by using this, we conclude that we have:
$$A\in M\implies A^c\in M$$

Next, let us look at unions. Assume $A_i\in M$, and consider their union:
$$A=\bigcup_{i=1}^\infty A_i$$

With $K\subset X$ being an arbitrary compact set, as above, set $B_1=A_1\cap K$, and then define recursively the following sets:
$$B_n=(A_n\cap K)-(B_1\cup\ldots\cup B_{n-1})$$

In terms of these sets $B_n$, we have the following formula:
$$A\cap K=\bigcup_{i=1}^\infty B_n$$

On the other hand, according to our definition of the sets $B_n$, and by using (12), we have $B_n\in M$. Thus, by using (10), we obtain $A\in M$. Thus, we have proved that:
$$A_i\in M\implies \bigcup_{i=1}^\infty A_i\in M$$

Finally, if a subset $C\subset X$ is closed, then $C\cap K$ is compact, and so $C\in M$. Thus, as claimed above, $M$ is indeed an algebra, which contains all the Borel sets.

\medskip

(14) Our next claim, which together with what we found in (10) will prove that $\mu$ is indeed a measure on $M$, is that, in the context of our constructions in (3), the family $N$ constructed there consists precisely of the sets $E\in M$ satisfying:
$$\mu(E)<\infty$$

In order to prove this, let $E\in N$. By using (7) and (12) we deduce that we have $E\cap K\in N$ for any compact set $K\subset X$, and we conclude from this that we have:
$$E\in N\implies E\in M$$

Conversely, assume that $E\in M$ has finite measure, $\mu(E)<\infty$. In order to prove $E\in N$, we pick a number $\varepsilon>0$. We can then find an open set $U\subset X$ such that:
$$E\subset U\quad,\quad\mu(U)<\infty$$

Also, by (8) and (11) we can find a compact set $K\subset X$ such that:
$$K\subset V\quad,\quad\mu(U-K)<\varepsilon$$

Now since $E\cap K\in N$, we can find a compact set $H\subset X$ such that:
$$H\subset E\cap K\quad,\quad\mu(E\cap K)<\mu(H)+\varepsilon$$

Observe now that we have the following inclusion:
$$E\subset(E\cap K)\cup(U-K)$$

By using this, we obtain the following inequality:
\begin{eqnarray*}
\mu(E)
&\leq&\mu\left((E\cap K)\cup(U-K)\right)\\
&\leq&\mu(E\cap K)+\mu(U-K)\\
&<&\mu(H)+2\varepsilon
\end{eqnarray*}

Thus $E\in N$, which finishes the proof of our claim, and proves that $\mu$ is a measure.

\medskip

(15) Time now to get into the final verification, that of the following formula:
$$I(f)=\int_Xf(x)d\mu(x)$$

As a first observation, by taking real and imaginary parts, it is enough to prove this for the real functions $f$. Moreover, and assuming now that $f$ is real, by using $f\to-f$, in order to prove the above equality, it is enough to prove the following inequality:
$$I(f)\leq\int_Xf(x)d\mu(x)$$

(16) So, let us prove now this latter inequality. For this purpose, let $K\subset X$ be the support of our function $f$, choose an interval containing the range, $Im(f)\subset[a,b]$, pick an arbitrary $\varepsilon>0$, and then numbers $y_0,\ldots,y_n$ such that $y_i-y_{i-1}<\varepsilon$, and:
$$y_0<a<y_1<\ldots<y_n=b$$

With this done, consider now the following sets:
$$E_i=\left\{x\in X\Big|y_{i-1}<f(x)<y_i\right\}\cap K$$

Since $f$ is continuous, these are disjoint Borel sets, whose union is $K$. Thus, we can find open sets $U_i\subset X$, subject to the following conditions:
$$E_i\subset U_i\quad,\quad \mu(U_i)<\mu(E_i)+\frac{\varepsilon}{n}\quad,\quad f(x)<y_i+\varepsilon,\forall x\in U_i$$

By using now Theorem 6.11, we can find certain continuous functions $h_i:X\to[0,1]$, with $h_i$ supported on $U_i$, such that the following equality holds, on $K$:
$$h_1(x)+\ldots+h_n(x)=1$$

By multiplying by $f$, we conclude that the following equality holds, on $K$:
$$f(x)=h_1(x)f(x)+\ldots+h_n(x)f(x)$$

On the other hand, by using (7), we have the following estimate:
$$\mu(K)\leq I\left(\sum_{i=1}^nh_i\right)=\sum_{i=1}^nI(h_i)$$

Good news, with all these ingredients in hand, we can now finish. Since we have $h_if\leq(y_i+\varepsilon)h_i$, and since $y_i-\varepsilon<f(x)$ on $E$, we have the following estimate:
\begin{eqnarray*}
I(f)
&=&\sum_{i=1}^nI(h_if)\\
&\leq&\sum_{i=1}^n(y_i+\varepsilon)I(h_i)\\
&=&\sum_{i=1}^n(|a|+y_i+\varepsilon)I(h_i)-|a|\sum_{i=1}^nI(h_i)\\
&\leq&\sum_{i=1}^n(|a|+y_i+\varepsilon)\left(\mu(E_i)+\frac{\varepsilon}{n}\right)-|a|\mu(K)\\
&=&\sum_{i=1}^n(y_i-\varepsilon)\mu(E_i)+2\varepsilon\mu(K)+\frac{\varepsilon}{n}\sum_{i=1}^n(|a|+y_i+\varepsilon)\\
&\leq&\int_Xf(x)d\mu(x)+\varepsilon\left(2\mu(K)+|a|+b+\varepsilon\right)
\end{eqnarray*}

Now since $\varepsilon>0$ was arbitrary, this finishes the proof of our latest claim, that is, of the inequality in (15), and so finishes the proof of the present theorem.
\end{proof}

\end{document}
